\documentclass[UTF8,a4paper,11pt]{ctexart}
\usepackage[a4paper,margin=2.3cm]{geometry}
\usepackage{amsmath,amssymb,mathtools}
\usepackage{enumitem}
\usepackage[hidelinks]{hyperref}
\usepackage{fancyhdr}
\setlist[itemize]{leftmargin=2.2em,itemsep=0.25em,topsep=0.35em}
\setlength{\parindent}{2em}\setlength{\parskip}{0.35em}
\pagestyle{fancy}\fancyhf{}\fancyhead[L]{离散数学与结构（I）}\fancyhead[R]{Lec 12}\fancyfoot[C]{\thepage}\setlength{\headheight}{14pt}
\title{组合计数}\author{zhangjunshuo}\date{2026 年 10 月}
\begin{document}
\maketitle
\begin{quote}
本讲根据 2026 年 10 月 22 日的两段课堂转写、三页手写原稿和课程教材整理。手写稿与转写都是部分覆盖：手写稿以清单、公式与推导为主（含格路/括号推导、广义二项式系数、线性递推特征方程、三角剖分与容斥的指示函数证明），转写补充口述解释、具体计数例子与部分细节；整理版按"取并集"处理，两边冲突或疑似笔误处给出勘误。

本讲主线：基本计数原理 $\to$ 二项式系数与组合恒等式 $\to$ 卡特兰数 $\to$ 球与桶 $\to$ 容斥原理与抽屉原理（Erdős--Szekeres 定理）。

每个分节开头都先交代本节用到的记号；所有符号在首次使用处给出定义，不依赖读者记得前面讲次的写法。
\end{quote}
\tableofcontents\clearpage

\section{基本计数原理、排列与组合}


组合计数研究的是"满足某些条件的对象有多少个"。这一节先把最基础的三种计数工具立起来：\textbf{加法原理}、\textbf{乘法原理}与\textbf{双射计数}，然后用它们数出排列、组合与多重组合数，并给出二项式系数的几条基本恒等式。后面几节（组合恒等式、卡特兰数、球与桶、容斥与抽屉）都建立在这一节的记号与语言上。

\textbf{本节用到的记号：} $[n]=\{1,2,\dots,n\}$；$|A|$ 为有限集 $A$ 的元素个数；$A\sqcup B$ 表示不交并；$A\times B$ 为笛卡尔积；$n!$ 为阶乘；$n^{\underline{k}}=n(n-1)\cdots(n-k+1)$ 为\textbf{下降阶乘}；$\binom nk$ 为二项式系数；$\operatorname{inj}([k],[n])$ 表示从 $[k]$ 到 $[n]$ 的单射全体。

\subsection{加法原理与乘法原理}


\textbf{加法原理。} 若 $A_1,\dots,A_t$ 是两两不交的有限集，则

\[
|A_1\sqcup A_2\sqcup\cdots\sqcup A_t|=|A_1|+|A_2|+\cdots+|A_t|.
\]

一句话：\textbf{分类计数，类与类不重不漏，就把各类数量相加。} 例如午饭在 20 种面食与 30 种米饭套餐中选一种，共有 $20+30=50$ 种。用之前必须确认"两类没有重叠"，否则要用后面的容斥原理。

\textbf{乘法原理。} 若 $A,B$ 是有限集，则

\[
|A\times B|=|A|\cdot |B|.
\]

更一般地，若一个对象由 $t$ 步依次选定，第 $i$ 步有 $a_i$ 种选法且各步选择互不影响，则总选择数为 $a_1a_2\cdots a_t$。例如先在 10 道素菜中选一道、再在 10 道荤菜中选一道，共有 $10\times 10=100$ 种组合。

\textbf{基本应用：函数计数。} 从 $[n]$ 到 $[m]$ 的函数共有 $m^n$ 个：$[n]$ 中每个元素独立地在 $m$ 个像中选一个。这是后面"球与桶"模型中"可区分球、可区分桶"的计数 $m^n$ 的来源。

\subsection{阶乘、排列数与单射计数}


\textbf{阶乘。} $n!=1\cdot2\cdots n$（约定 $0!=1$）。它等于 $[n]$ 的全部排列数，也就是 $[n]$ 到自身的双射个数：选第一个位置有 $n$ 种，第二个 $n-1$ 种，依次相乘。

\textbf{排列数（下降阶乘）。} 从 $n$ 个不同对象中有顺序地选出 $k$ 个，方案数为

\[
P(n,k)=n(n-1)\cdots(n-k+1)=n^{\underline{k}}=\frac{n!}{(n-k)!}\qquad(0\le k\le n).
\]

\textbf{组合意义。} 它恰好是\textbf{单射} $f:[k]\to[n]$ 的个数：第一个位置（输入 $1$）选一个像有 $n$ 种，第二个位置（输入 $2$）不能与前面重复，剩 $n-1$ 种，依此类推。所以

\[
\#\operatorname{inj}([k],[n])=n^{\underline{k}}=\frac{n!}{(n-k)!}.
\]

例如 $P(5,3)=5\cdot4\cdot3=60$。当 $k>n$ 时约定 $n^{\underline{k}}=0$（放不下），于是公式对所有 $k\ge0$ 都成立。

\subsection{组合数与二项式系数}


\textbf{定义。} 从 $n$ 个不同对象中\textbf{不重复、不计顺序}地选出 $k$ 个，方案数记作

\[
\binom nk=\frac{n^{\underline{k}}}{k!}=\frac{n!}{k!\,(n-k)!}\qquad(0\le k\le n),
\]

称为\textbf{二项式系数}（也记 $C(n,k)$、$C_n^k$）。它计数的是 $[n]$ 的 $k$ 元子集：先按顺序选（$n^{\underline{k}}$ 种），再除掉 $k$ 个元素的 $k!$ 种排列。于是

\[
P(n,k)=\binom nk\,k! .
\]

约定：$k<0$ 或 $k>n$ 时 $\binom nk=0$；$\binom n0=\binom nn=1$。

\textbf{组合数与杨辉三角。} 把 $\binom nk$ 按 $n$ 排成三角形，就是帕斯卡三角：每一行两端是 $1$，中间每个数等于上一行左右两个数之和——这正是下面的帕斯卡恒等式。

\subsection{二项式系数的四条基本恒等式}


\textbf{恒等式一（对称性）。}

\[
\binom nk=\binom n{n-k}.
\]

\textbf{证明（双射）。} $[n]$ 的 $k$ 元子集 $S$ 与其余集 $[n]\setminus S$（$n-k$ 元子集）一一对应。$\square$

\textbf{恒等式二（帕斯卡恒等式）。} 对 $1\le k\le n-1$，

\[
\binom nk=\binom{n-1}{k-1}+\binom{n-1}{k}.
\]

\textbf{证明（分类）。} 看 $[n]$ 的 $k$ 元子集是否包含元素 $n$：包含 $n$ 的，相当于在 $[n-1]$ 中再选 $k-1$ 个，共 $\binom{n-1}{k-1}$ 种；不包含 $n$ 的，相当于在 $[n-1]$ 中选 $k$ 个，共 $\binom{n-1}{k}$ 种。两类不交且覆盖全部，由加法原理得证。$\square$

\textbf{恒等式三（吸收恒等式）。} 对 $1\le k\le n$，

\[
k\binom nk=n\binom{n-1}{k-1}.
\]

\textbf{证明（双计数）。} 左边数的是"先选 $k$ 元子集、再从其中指定一个元素"的方案数：先选子集 $\binom nk$ 种，再在 $k$ 个元素中选一个 $k$ 种。右边数的是"先指定一个元素、再从剩下的 $n-1$ 个里选 $k-1$ 个与它组成子集"：指定元素 $n$ 种，再选其余 $k-1$ 个 $\binom{n-1}{k-1}$ 种。两边数的是同一批对象（一个带标记的 $k$ 元子集）。$\square$

代数写法：$\binom nk=\frac nk\binom{n-1}{k-1}$。这条恒等式是把"$k$ 的幂"化成二项式系数的钥匙，后面求 $\sum k\binom nk$ 时反复使用。

\textbf{恒等式四（乘积恒等式）。} 对 $0\le k\le m\le n$，

\[
\binom nm\binom mk=\binom nk\binom{n-k}{m-k}.
\]

\textbf{证明（双计数）。} 左边：先从 $[n]$ 选一个 $m$ 元子集，再从中选一个 $k$ 元子集，等价于把 $[n]$ 分成三组（大小 $k$、$m-k$、$n-m$），方案数 $\frac{n!}{k!(m-k)!(n-m)!}$。右边：先选 $k$ 元子集，再从剩下的 $n-k$ 个中选 $m-k$ 个，同样得到同一批"三组划分"，方案数相同。$\square$

\textbf{组合证明的两种写法。} 上面反复用了两种方法：\textbf{(1) 双射}：把待数对象与已知对象一一对应；\textbf{(2) 双计数}：把同一个集合用两种方式数出来，两个结果必须相等。做组合题时先问"这两边分别数的是什么集合"，往往比代数变形更快。

\subsection{多重组合数与多项式展开的系数}


\textbf{定义（多重组合数）。} 设 $k_1,\dots,k_t\ge0$、$k_1+\cdots+k_t=n$。把 $n$ 个不同对象分成\textbf{有序的} $t$ 组、第 $i$ 组恰有 $k_i$ 个的方案数为

\[
\binom{n}{k_1,\dots,k_t}=\frac{n!}{k_1!\,k_2!\cdots k_t!}.
\]

\textbf{证明。} 依次从剩余对象中选：先选第 1 组 $\binom n{k_1}$ 种，再从剩下的 $n-k_1$ 个中选第 2 组 $\binom{n-k_1}{k_2}$ 种，依此类推，乘积化简即得。特别地，$t=2$ 时 $\binom{n}{k,n-k}=\binom nk$。$\square$

\textbf{二项式定理。} 对任意 $x,y$ 与 $n\ge0$，

\[
(x+y)^n=\sum_{k=0}^{n}\binom nk x^k y^{n-k}.
\]

\textbf{证明（组合意义）。} 把 $(x+y)^n$ 看成 $n$ 个因子 $(x+y)$ 相乘。展开时每个因子贡献一个 $x$ 或一个 $y$；若恰有 $k$ 个因子贡献 $x$，则得到单项式 $x^ky^{n-k}$。选择哪 $k$ 个因子贡献 $x$ 有 $\binom nk$ 种，故该项系数为 $\binom nk$。$\square$

（也可以用帕斯卡恒等式对 $n$ 归纳证明：从 $n-1$ 到 $n$ 时，$x^ky^{n-k}$ 的系数由上一行的 $\binom{n-1}{k-1}$ 与 $\binom{n-1}{k}$ 相加得到。）

\textbf{多项式定理。} 类似地，

\[
(x_1+x_2+\cdots+x_t)^n=\sum_{\substack{k_1,\dots,k_t\ge0\\ k_1+\cdots+k_t=n}}\binom{n}{k_1,\dots,k_t}x_1^{k_1}x_2^{k_2}\cdots x_t^{k_t}.
\]

例如 $(x+y+z)^n=\sum_{a+b+c=n}\frac{n!}{a!b!c!}x^ay^bz^c$。

\textbf{两个常用特值。} 在二项式定理中取 $x=y=1$ 得

\[
\sum_{k=0}^{n}\binom nk=2^n,
\]

即 $[n]$ 的全部子集个数；取 $x=1,y=-1$（$n\ge1$）得

\[
\sum_{k=0}^{n}(-1)^k\binom nk=0,
\]

即偶数元子集与奇数元子集一样多。

\subsection{本节的勘误与提醒}


\begin{itemize}
\item \textbf{加法原理要求"不交"。} 若两类有重叠，直接相加会重复计数，应改用容斥原理。
\item \textbf{排列数 vs 组合数。} 计顺序用 $P(n,k)$（单射），不计顺序用 $\binom nk$（子集）；两者差一个 $k!$。
\item \textbf{$\binom nk$ 的约定要统一。} 本节取 $k<0$ 或 $k>n$ 时为 $0$（$n\ge0$ 时）；第 2 节把 $n$ 推广到负数时，$k>n$ 不再自动为 $0$，要按下降阶乘的定义算。
\item \textbf{多重组合数是有序分组。} $\binom{n}{k_1,\dots,k_t}$ 中组的编号是固定的；若组不标号，要再除以组的对称性，那属于另一类问题（分拆数）。
\item 课堂转写把"下降阶乘/排列数"的记号念成"$n$ 的 $k$ 次方"或"$A_n^k$"，以本节的 $n^{\underline{k}}=\frac{n!}{(n-k)!}$ 为准。
\end{itemize}

\section{组合恒等式与求和技巧}


上一节给出了二项式系数与二项式定理。这一节把它们变成"可运算的"工具：先证明几条最常用的求和恒等式（曲棍球恒等式、Vandermonde、加权和），再说明"把 $k$ 的多项式拆成二项式系数"这一机械化方法，最后把二项式系数推广到负的 $n$，并用一个例子（$A_m=\sum_k(-1)^k\binom{m-k}{k}$）引出线性递推的特征方程解法。

\textbf{本节用到的记号：} $\binom nk$ 为二项式系数；约定 $\binom nk=0$（$k<0$ 或 $n\ge0$ 且 $k>n$）；$n^{\underline{k}}$ 为下降阶乘；对固定 $j$，$\binom kj$ 视为 $k$ 的函数；$\lambda$ 为递推的特征根。

\subsection{曲棍球恒等式（斜向求和）}


\textbf{定理一。} 对 $0\le m\le n$，

\[
\sum_{k=m}^{n}\binom km=\binom{n+1}{m+1}.
\]

\textbf{证明（双计数）。} 右边数的是 $[n+1]$ 的 $(m+1)$ 元子集。对每个这样的子集，看它的\textbf{最大元素}：若最大元素为 $k+1$（其中 $k$ 从 $m$ 到 $n$），则其余 $m$ 个元素从 $[k]=\{1,\dots,k\}$ 中选，共 $\binom km$ 种。按最大元素分类求和即得左边。$\square$

这条恒等式在帕斯卡三角上表现为一条斜线的和，因此叫\textbf{曲棍球恒等式}（hockey-stick identity）。特例 $m=1$：$\sum_{k=1}^n k=\binom{n+1}{2}$；$m=2$：$\sum_{k=2}^n\binom k2=\binom{n+1}{3}$。

\textbf{定理二（另一种斜向求和，一般 $r$ 版）。} 对整数 $r\ge0$、$n\ge0$，

\[
\sum_{k=0}^{n}\binom{r+k}{k}=\binom{r+n+1}{n}.
\]

\textbf{证明（双计数）。} 右边数 $[r+n+1]$ 的 $n$ 元子集 $S$。因为 $|S|=n<r+n+1$，$S$ 之外必有元素，设 $j$ 是\textbf{最小的不属于 $S$ 的元素}。令 $k=n+1-j$，则 $k$ 取值 $0,1,\dots,n$；此时 $1,2,\dots,j-1=n-k$ 都在 $S$ 中，$j$ 不在 $S$ 中，而 $S$ 还要在 $\{j+1,\dots,r+n+1\}$ 中选剩下的 $k$ 个元素。该集合的大小是

\[
(r+n+1)-(n-k+1)=r+k,
\]

所以对应 $\binom{r+k}{k}$ 个 $S$。按 $j$（也就是按 $k$）分类求和即得。取 $r=n$ 即得特例

\[
\sum_{k=0}^{n}\binom{n+k}{k}=\binom{2n+1}{n}.
\]

$\square$

\textbf{例。} 取 $r=3,n=3$：左边 $=\binom30+\binom41+\binom52+\binom63=1+4+10+20=35$，右边 $\binom{7}{3}=35$。

\subsection{Vandermonde 恒等式}


\textbf{定理三（Vandermonde）。} 对 $r,m,n\ge0$，

\[
\sum_{k=0}^{n}\binom rk\binom m{n-k}=\binom{r+m}{n}.
\]

\textbf{证明（双计数）。} 右边数 $r+m$ 个对象中选 $n$ 个。把对象分成两组：$r$ 个与 $m$ 个。设选出的 $n$ 个中恰有 $k$ 个来自第一组，则来自第二组的有 $n-k$ 个；按 $k$ 分类求和即得。$\square$

\textbf{推论（平方和）。} 取 $r=m=n$，并用对称性 $\binom m{n-k}=\binom nk$，得

\[
\sum_{k=0}^{n}\binom nk^2=\binom{2n}{n}.
\]

\textbf{例。} $n=3$：$\binom30^2+\binom31^2+\binom32^2+\binom33^2=1+9+9+1=20=\binom63$。

\subsection{加权和：把 $k$ 降成组合数}


求 $\sum_k k\binom nk$、$\sum_k k^2\binom nk$ 这类和式，核心是把"$k$ 的幂"改写成"组合数"，再利用下面的双计数恒等式。

\textbf{定理四（核心引理）。} 对 $0\le j\le n$，

\[
\sum_{k=0}^{n}\binom kj\binom nk=\binom nj 2^{n-j}.
\]

\textbf{证明（双计数）。} 左边数的是"先选 $k$ 元子集、再从其中选 $j$ 元子集"，等价于"先选 $j$ 元子集，再在剩下的 $n-j$ 个元素中任选若干个"：先选 $\binom nj$ 种，剩下的 $n-j$ 个每个都有选/不选两种，共 $2^{n-j}$ 种。$\square$

\textbf{推论一。} 取 $j=1$：因为 $\binom k1=k$，

\[
\sum_{k=0}^{n}k\binom nk=\binom n1 2^{n-1}=n\,2^{n-1}.
\]

\textbf{推论二。} 取 $j=2$，并用 $k^2=k(k-1)+k=2\binom k2+\binom k1$：

\[
\sum_{k=0}^{n}k^2\binom nk
=2\sum_k\binom k2\binom nk+\sum_k\binom k1\binom nk
=2\binom n2 2^{n-2}+n\,2^{n-1}
=n(n+1)2^{n-2}.
\]

\textbf{推论三（平方和公式的另一种证法）。} 对普通平方和，取 $\binom nk$ 换成 $1$ 不适用；改用 $k^2=2\binom k2+\binom k1$ 与曲棍球恒等式：

\[
\sum_{k=0}^{n}k^2=2\binom{n+1}{3}+\binom{n+1}{2}=\frac{n(n+1)(2n+1)}{6}.
\]

\textbf{机械化方法。} 任何关于 $k$ 的多项式都能唯一写成 $\binom kj$ 的线性组合（等价地，写成下降阶乘之和）；再套定理四即可。例：$k^3=6\binom k3+6\binom k2+\binom k1$，于是

\[
\sum_k k^3\binom nk=6\binom n3 2^{n-3}+6\binom n2 2^{n-2}+n2^{n-1}.
\]

这就是课堂上说的"对任何一个多项式，都有一种比较机械化的方式先把它变成若干二项式系数之和"。

\subsection{广义二项式系数：让 $n$ 取负整数}


\textbf{定义。} 对任意整数 $n$（包括负数）与 $k\ge0$，定义

\[
\binom nk=\frac{n(n-1)\cdots(n-k+1)}{k!}=\frac{n^{\underline{k}}}{k!},
\]

并约定 $k<0$ 时 $\binom nk=0$。当 $n\ge0$、$k>n$ 时，乘积中会出现因子 $0$，所以仍然是 $0$；但 $n<0$ 时不会自动为零。

\textbf{例。} $\binom{-1}{k}=\frac{(-1)(-2)\cdots(-k)}{k!}=(-1)^k$。更一般地，

\[
\binom{-m}{k}=(-1)^k\binom{m+k-1}{k}\qquad(m\ge1).
\]

\textbf{为什么有用。} 帕斯卡恒等式与吸收恒等式对任意整数 $n$ 仍然成立：

\[
\binom nk=\binom{n-1}{k}+\binom{n-1}{k-1},\qquad k\binom nk=n\binom{n-1}{k-1}.
\]

于是很多"对正整数 $n$ 证明的组合恒等式"可以自然地延拓到负数；反过来，对负数 $n$ 的恒等式常能化成对正整数的新恒等式。它与形式幂级数

\[
(1+x)^{-m}=\sum_{k\ge0}\binom{-m}{k}x^k=\sum_{k\ge0}(-1)^k\binom{m+k-1}{k}x^k
\]

配合，是生成函数与线性递推的标准工具（本课程下一章生成函数会系统使用）。

\subsection{例子：$A_m=\sum_k(-1)^k\binom{m-k}{k}$ 与线性递推}


\textbf{定义。} 令

\[
A_m=\sum_{k=0}^{\lfloor m/2\rfloor}(-1)^k\binom{m-k}{k}\qquad(m\ge0).
\]

这里上标 $m-k$ 始终非负，所以 $\binom{m-k}{k}$ 就是第 1 节的普通二项式系数，与第 2.4 小节"$n$ 取负整数"的推广无关。（等价写法是 $\sum_{k\ge0}$ 并规定 $k>m-k$ 时该项为 $0$。）

\textbf{递推。} 由帕斯卡恒等式

\[
\binom{m-k}{k}=\binom{m-1-k}{k}+\binom{m-1-k}{k-1},
\]

代入并对 $k$ 求和：

\[
A_m=\sum_{k\ge0}(-1)^k\binom{m-1-k}{k}+\sum_{k\ge0}(-1)^k\binom{m-1-k}{k-1}
=A_{m-1}-A_{m-2},
\]

第二个等号中令 $j=k-1$ 并用 $(-1)^{j+1}=-(-1)^j$。这里出现的 $\binom{m-1-k}{k}$ 等仍按第 1 节的\textbf{普通约定}理解（$k$ 超过上标时为 $0$），不套用第 2.4 小节负 $n$ 的定义；递推从 $m\ge2$ 起使用。初值：$A_0=\binom00=1$，$A_1=\binom10-\binom01=1$。于是

\[
A_0,A_1,A_2,A_3,A_4,A_5,A_6,\dots=1,1,0,-1,-1,0,1,\dots
\]

\textbf{解递推（特征方程法）。} 对二阶线性递推

\[
a_m=\alpha a_{m-1}+\beta a_{m-2},
\]

设 $a_m=\lambda^m$ 代入得

\[
\lambda^2=\alpha\lambda+\beta .
\]

\begin{itemize}
\item 若有两个相异根 $\lambda_1\neq\lambda_2$，则通解为 $a_m=C_1\lambda_1^m+C_2\lambda_2^m$；
\item 若有重根 $\lambda$，则通解为 $a_m=(C_1+C_2m)\lambda^m$。
\end{itemize}


常数 $C_1,C_2$ 由初值确定。（验证：两种形式代入递推都成立；两个参数恰好匹配两个初值。）

对 $A_m$：方程 $\lambda^2=\lambda-1$，即 $\lambda^2-\lambda+1=0$，两根为 $\lambda=e^{\pm i\pi/3}$（六次本原单位根）。因此

\[
A_m=C\,e^{im\pi/3}+\overline{C}\,e^{-im\pi/3}.
\]

由 $A_0=1$、$A_1=1$ 解得 $C=\frac12-\frac{i}{2\sqrt3}$，即

\[
A_m=\cos\frac{m\pi}{3}+\frac1{\sqrt3}\sin\frac{m\pi}{3}
=\frac{2}{\sqrt3}\sin\frac{(m+1)\pi}{3}.
\]

特别地 $A_m$ 以 $6$ 为周期。

\textbf{这个例子的意义。} 它说明"带 $(-1)^k$ 的组合和"可以化成线性递推，从而用特征方程求出闭式——这是组合恒等式与递推方法之间的桥梁。

\subsection{本节的勘误与提醒}


\begin{itemize}
\item \textbf{手写稿底部的 $A_m$ 与 $A_n$ 是两种不同的和，不要读混。} 手写稿（高倍放大可辨）写的是
\end{itemize}


\[
  \begin{aligned}
  A_m&=\sum_{k\ge0}\binom{k}{m}2^{-k}
  =\sum_k\binom{k-1}{m}2^{-k}+\sum_k\binom{k-1}{m-1}2^{-k}\\
  &=\sum_k\binom{k}{m}2^{-k-1}+\sum_k\binom{k}{m-1}2^{-k-1}
  =\frac12A_m+\frac12A_{m-1}.
  \end{aligned}
\]

  下标是 $m$、上标是 $k$。这个推导是\textbf{正确}的：由 $\sum_{k\ge m}\binom{k}{m}x^k=\frac{x^m}{(1-x)^{m+1}}$ 取 $x=1/2$ 得 $A_m=2$，所以 $A_m=\frac12A_m+\frac12A_{m-1}$ 与 $A_m=A_{m-1}$ 都成立（两边都等于 $2$）。它是"换指标 $+$ 折半"技巧的演示，不是错误。

  但注意另一种容易混淆的和：若定义 $A_n=\sum_{k=0}^{n}\binom{n}{k}2^{-k}$（\textbf{上标是 $n$、下标是 $k$}，与上面的 $A_m$ 不同），则由二项式定理 $A_n=(3/2)^n$，正确的递推是 $A_n=A_{n-1}+\frac12A_{n-1}=\frac32A_{n-1}$；此时末行 $\frac12A_n+\frac12A_{n-1}$ 不成立（反例 $n=2$：左边 $9/4$，右边 $15/8$）。做题时先看清二项式系数是"上标变化"还是"下标变化"。
\begin{itemize}
\item \textbf{手写稿 $\sum_k k^2$ 一行的小笔误。} 该行把 $2\binom{n+1}{3}$ 写成了 $2\binom{k+1}{3}$，上标应为 $n+1$；最终结果 $\frac{n(n+1)(2n+1)}6$ 正确。
\item \textbf{曲棍球恒等式的方向。} 是 $\sum_{k=m}^n\binom km=\binom{n+1}{m+1}$，不是 $\binom{n+1}{m}$；注意"上限 $+1$、$m$ 也 $+1$"。
\item \textbf{定理四要求 $j\le n$。} 当 $j>n$ 时两边都是 $0$（按约定），公式仍成立；但若把 $\binom kj$ 换成别的函数，机械化方法不一定成立。
\item \textbf{特征方程法只适用于常系数线性递推。} 系数随 $n$ 变化时（如 $a_n=na_{n-1}$）不能用 $\lambda^n$ 试解，那要靠生成函数或其它方法。
\item \textbf{重根时通解是 $(C_1+C_2m)\lambda^m$，不是 $C\lambda^m$。} 课堂转写中"$A_n=n\lambda^n$"的说法是重根情形的简化形式，需补上另一项。
\item 转写里"广义二项式系数 $\binom{-1}{k}$"一段识别得很乱；以本节定义与 $\binom{-1}{k}=(-1)^k$ 为准。
\end{itemize}

\section{卡特兰数}


\textbf{卡特兰数}（Catalan numbers）是组合计数里出现频率最高的一列数：括号匹配、格路、多边形三角剖分、二叉树、加括号方式，数出来都是同一列 $1,1,2,5,14,42,\dots$。这一节先给出它的定义与递推，再用格路与反射原理证明通项公式 $C_n=\frac1{n+1}\binom{2n}{n}$，最后说明它与三角剖分、加括号的对应关系。

\textbf{本节用到的记号：} $C_n$ 为第 $n$ 个卡特兰数（约定从 $C_0=1$ 开始）；格路指每次向右 $(1,0)$ 或向上 $(0,1)$ 走一步的路径；$\binom nk$ 为二项式系数。

\subsection{从括号序列说起}


考虑 $n$ 对左右括号组成的序列。称它是\textbf{合法的}（balanced），若从左往右读时，任意前缀中左括号的个数都不少于右括号的个数；等价地，整个序列左右括号各 $n$ 个且前缀平衡条件成立。

\begin{itemize}
\item $n=1$：只有 \texttt{()} 合法，共 $1$ 个；
\item $n=2$：\texttt{(())}、\texttt{()()} 合法，而 \texttt{())(}、\texttt{)(()} 等不合法，共 $2$ 个；
\item $n=3$：合法序列有 \texttt{((()))}、\texttt{(()())}、\texttt{(())()}、\texttt{()(())}、\texttt{()()()}，共 $5$ 个。
\end{itemize}


记 $C_n$ 为 $n$ 对括号的合法序列数。直接计算可得前几项

\[
C_0,C_1,C_2,C_3,C_4,C_5,C_6=1,1,2,5,14,42,132.
\]

\textbf{递推关系。} 任何 $n\ge1$ 的合法序列都以左括号开头。设第一个左括号所匹配的右括号在第 $2k+2$ 个位置（$k=0,1,\dots,n-1$）。那么：

\begin{itemize}
\item 第一个左括号与它匹配的右括号之间的 $2k$ 个括号，本身是一个 $k$ 对括号的合法序列，有 $C_k$ 种；
\item 这个右括号之后的剩余部分，是一个 $n-1-k$ 对括号的合法序列，有 $C_{n-1-k}$ 种。
\end{itemize}


按 $k$ 分类相加：

\[
C_n=\sum_{k=0}^{n-1}C_k\,C_{n-1-k}\qquad(n\ge1),\qquad C_0=1.
\]

这个递推是卡特兰数最本质的定义；后面几种模型都会回到它。

\subsection{格路与反射原理}


\textbf{格路模型。} 从 $(0,0)$ 出发，每次向右 $(1,0)$ 或向上 $(0,1)$，走到 $(n,n)$，且要求路径\textbf{不越过对角线} $y=x$（即始终满足 $y\le x$）。这样的路径称为 Dyck 路径。

\textbf{与括号的对应。} 把"向右"读成左括号、"向上"读成右括号：走 $n$ 步向右、$n$ 步向上，对应 $n$ 对括号；$y\le x$ 等价于"左括号数不少于右括号数"。因此 Dyck 路径数也是 $C_n$。

\textbf{总路径数。} 从 $(0,0)$ 到 $(n,n)$ 共走 $2n$ 步，其中选 $n$ 步向右，故总路径数为 $\binom{2n}{n}$。

\textbf{坏路径计数（反射原理）。} 称越过对角线的路径为坏路径，即它碰到直线 $y=x+1$。对每条坏路径，取它\textbf{第一次}碰到 $y=x+1$ 的时刻，把从起点到该时刻的这一段沿直线 $y=x+1$ 反射。起点 $(0,0)$ 反射后变成 $(-1,1)$，而反射不改变步数，于是得到一条从 $(-1,1)$ 到 $(n,n)$ 的路径；反过来也唯一还原。故坏路径与"从 $(-1,1)$ 到 $(n,n)$ 的所有路径"一一对应。

从 $(-1,1)$ 到 $(n,n)$ 需要向右 $n+1$ 步、向上 $n-1$ 步，共 $2n$ 步，因此路径数为

\[
\binom{2n}{n+1}=\binom{2n}{n-1}.
\]

于是

\[
C_n=\binom{2n}{n}-\binom{2n}{n-1}
=\binom{2n}{n}\left(1-\frac{n}{n+1}\right)
=\frac{1}{n+1}\binom{2n}{n}.
\]

这就是卡特兰数的通项公式。

\textbf{例（$n=4$）。} 总路径数 $\binom84=70$，坏路径数 $\binom83=56$，好路径数 $70-56=14=C_4$。

\subsection{多边形三角剖分与加括号}


\textbf{三角剖分。} 一个凸 $(n+2)$ 边形，用 $n-1$ 条互不相交的对角线把它剖分成 $n$ 个三角形，称为一个\textbf{三角剖分}。三角剖分的个数也是 $C_n$。

\textbf{证明（按含固定边的三角形分类）。} 设多边形顶点依次为 $v_1,v_2,\dots,v_{n+2}$，固定边 $v_1v_2$。在任何三角剖分中，边 $v_1v_2$ 恰好属于某一个三角形，设这个三角形的第三个顶点为 $v_k$（$3\le k\le n+2$）。这条对角线把多边形分成两个较小的多边形：

\begin{itemize}
\item 顶点为 $v_2,v_3,\dots,v_k$ 的 $(k-1)$ 边形；
\item 顶点为 $v_1,v_k,v_{k+1},\dots,v_{n+2}$ 的 $(n+4-k)$ 边形。
\end{itemize}


前者的三角剖分数是 $C_{k-3}$（因为 $(k-1)$ 边形对应下标 $k-3$），后者是 $C_{n+2-k}$（$(n+4-k)$ 边形对应下标 $n+2-k$）。令 $a=k-3$，则 $a$ 从 $0$ 到 $n-1$，得到

\[
C_n=\sum_{a=0}^{n-1}C_a\,C_{n-1-a},
\]

正是卡特兰数的递推。$\square$

\textbf{加括号。} 给 $n+1$ 个符号的乘积加\textbf{完全括号}，不同的加括号方式数也是 $C_n$。例如 $n=2$（3 个符号 $a,b,c$）有 $(ab)c$ 与 $a(bc)$ 两种。

\textbf{为什么与三角剖分一致。} 固定 $(n+2)$ 边形的一条边 $e$ 作为"最终结果"，其余 $n+1$ 条边各看成一个符号。三角剖分中每个三角形看成一次乘法：它把三条边对应的对象合并成一个新对象；从叶子三角形开始逐层合并，最后在 $e$ 上得到整个乘积。不同的剖分给出不同的合并顺序，也就给出不同的加括号方式。反过来，给定 $n+1$ 个符号的完全加括号方式，从最内层开始把每次乘法画成一个三角形、逐层向外扩展，就唯一地得到 $(n+2)$ 边形的一个三角剖分。因此

\[
\text{三角剖分}\ \longleftrightarrow\ \text{加括号方式}\ \longleftrightarrow\ \text{合法括号序列},
\]

三者数量都等于 $C_n$。

\subsection{其他等价模型与提醒}


\textbf{等价模型。} $C_n$ 还等于：

\begin{itemize}
\item $n$ 个\textbf{内部结点}（$n+1$ 片叶子）的满二叉树的个数；
\item $n$ 对括号的合法序列数；
\item 从 $(0,0)$ 到 $(n,n)$ 不越过对角线的格路数；
\item $(n+2)$ 边形的三角剖分数；
\item $n+1$ 个符号的完全加括号方式数；
\item $n$ 个元素经一个栈得到的\textbf{合法出栈序列}数（补充，见 [DS] §11.1 例题）。
\end{itemize}


\textbf{下标不要错位。} $C_n$ 对应"$n$ 对括号 / $n+2$ 边形 / $n+1$ 个符号 / 不越过对角线的 $n\times n$ 格路"。做题时先确认题目里的 $n$ 是哪一个，再套公式。

\textbf{公式自检。} $n=0$ 时 $\frac1{1}\binom00=1=C_0$；$n=1$ 时 $\frac12\binom21=1$；$n=2$ 时 $\frac13\binom42=2$。用递推算出的前几项与公式逐项一致。

\textbf{本节的勘误与提醒。}
\begin{itemize}
\item 手写稿里"$C_{2n}=\binom{2n}{n}-\binom{2n}{n-1}$"中的 $C_{2n}$ 是笔误，左边应是第 $n$ 个卡特兰数 $C_n$；$\binom{2n}{n}-\binom{2n}{n-1}$ 是\textbf{计数结果}，不是卡特兰数的下标。
\item 反射原理中"坏路径"要按\textbf{第一次}碰到 $y=x+1$ 反射，且反射的是起点到该时刻的一段；反射后终点不变，起点变成 $(-1,1)$。
\item 括号序列的合法条件是"任意前缀左括号不少于右括号"，不是"左右括号总数相等"就够。
\item 三角剖分中固定的是\textbf{一条边}（而不是一个顶点）；分类用的三角形必须含这条边，分类才不重不漏。
\item 课堂转写把"多边形三角剖分"和"加括号"分成两段讲，且第二段开头从多边形讲起；手写稿则先画格路。两处说的是同一件事，整理时按"格路 → 递推 → 三角剖分/加括号"的顺序合并。
\end{itemize}

\section{球与桶：四类分配模型}


"把 $n$ 个球放进 $m$ 个桶"是组合计数里最常用的模型。球是否可区分、桶是否可区分、是否允许空桶，这三个问题决定了四个不同的公式。这一节把课堂上讲到的几类逐一列清，并说明它们分别对应什么计数对象。

\textbf{本节用到的记号：} $n$ 为球数，$m$ 为桶数；"可区分"指每个球（或桶）有编号、彼此不同；"相同"指球（或桶）不可区分；$\binom nk$ 为二项式系数；$S(n,m)$ 为第二类 Stirling 数（把 $n$ 个不同元素分成 $m$ 个非空无标号组的方案数）。

\subsection{四个基本模型}


\textbf{模型 A：球可区分，桶可区分，允许空桶。}

每个球独立地在 $m$ 个桶中选一个，所以方案数为

\[
m^n.
\]

等价说法：从 $[n]$ 到 $[m]$ 的函数个数（把球编号为 $1,\dots,n$，桶编号为 $1,\dots,m$，函数值就是球所在的桶）。例：3 个不同的球放入 2 个不同的桶，共 $2^3=8$ 种。

\textbf{模型 B：球可区分，桶可区分，不允许空桶。}

这是\textbf{满射}计数（每个桶至少有一个球，即函数是满射）。它的公式由容斥原理得到（见第 5 节）：

\[
\#\{\text{满射 }[n]\to[m]\}=\sum_{i=0}^{m}(-1)^i\binom mi(m-i)^n .
\]

记作 $\operatorname{onto}(n,m)$。它与第二类 Stirling 数的关系是 $\operatorname{onto}(n,m)=m!\,S(n,m)$：先做无标号的分组（$S(n,m)$），再给 $m$ 个组分配桶的编号（$m!$ 种）。例：3 个不同的球放入 2 个不同的桶且每桶非空，共 $2^3-2=6$ 种。

\textbf{模型 C：球相同，桶可区分，不允许空桶。}

把 $n$ 个相同的球排成一列，球与球之间有 $n-1$ 个空隙，选 $m-1$ 个空隙插入隔板，就把它分成 $m$ 段、每段非空。方案数为

\[
\binom{n-1}{m-1}.
\]

等价说法：正整数解 $x_1+\cdots+x_m=n$ 的个数。例：10 个相同的球放入 4 个不同的桶且每桶非空，共 $\binom93=84$ 种。

\textbf{模型 D：球相同，桶可区分，允许空桶。}

两种看法：

\begin{itemize}
\item 先用"补球"把空桶问题化归为非空：若允许空桶，先给每个桶放 1 个球（共用 $m$ 个），再按模型 C 分配剩下的；等价地，$n$ 个相同球放入 $m$ 个桶、允许空 = $n+m$ 个相同球放入 $m$ 个桶、不允许空。方案数为
\end{itemize}


\[
\binom{n+m-1}{m-1}.
\]

\begin{itemize}
\item 隔板法：允许空段时，等价于在 $n$ 个球与 $m-1$ 个隔板共 $n+m-1$ 个位置中选 $m-1$ 个放隔板。
\end{itemize}


等价说法：非负整数解 $x_1+\cdots+x_m=n$ 的个数。例：10 个相同的球放入 4 个不同的桶，允许空桶，共 $\binom{13}{3}=286$ 种。

\textbf{汇总。} 四类模型可以这样记：

\begin{itemize}
\item 球可区分 + 桶可区分 + 允许空：$m^n$（每个球独立选桶）；
\item 球可区分 + 桶可区分 + 不允许空：满射数 $\sum_i(-1)^i\binom mi(m-i)^n$（容斥，见第 5 节）；
\item 球相同 + 桶可区分 + 不允许空：$\binom{n-1}{m-1}$（正整数解）；
\item 球相同 + 桶可区分 + 允许空：$\binom{n+m-1}{m-1}$（非负整数解）。
\end{itemize}


\textbf{边界情形。} 当 $n<m$ 时，不可能每桶非空，所以模型 B、C 的答案为 $0$（$\binom{n-1}{m-1}$ 按约定也为 $0$）；模型 A、D 仍有意义。

\subsection{两种"都不区分"的情形（补充）}


如果连桶也不区分，问题就变成"把 $n$ 分拆成至多 $m$ 个部分"或"把 $n$ 个元素分成 $m$ 组"，一般没有简单的闭式，需要专门的工具：

\begin{itemize}
\item \textbf{球相同、桶相同}：$n$ 的\textbf{整数分拆}数，记 $p(n,m)$ 为"把 $n$ 写成\textbf{恰好} $m$ 个正整数之和"的方案数（不计次序）。它没有闭式公式，但有递推 $p(n,m)=p(n-1,m-1)+p(n-m,m)$（分类：含 $1$ 的分拆去掉一个 $1$；不含 $1$ 的分拆每个部分减 $1$）。若问"至多 $m$ 个部分"，则等于 $p(n,1)+p(n,2)+\cdots+p(n,m)$。
\item \textbf{球可区分、桶相同}：第二类 Stirling 数 $S(n,m)$，即把 $n$ 个不同元素分成 $m$ 个非空无标号组的方案数；满射数除以 $m!$ 就是它。
\end{itemize}


课堂上主要用前四个模型（桶可区分），这两种"桶相同"的情形属于教材补充（[DS] §11.1、讲义 §9.5 的表），知道对应关系即可，本讲不要求计算。更完整的"十二重计数表"（对象/盒子是否可区分、是否允许空、是否要求满射等各种组合）见 [DS] 表 11.1 与讲义 §9.5；本讲只要求前四类。

\subsection{怎么判断用哪个公式}


拿到"分配"问题，按顺序问三个问题：

1. \textbf{球是否可区分？} 若每个球有编号（不同的人、不同的球），用"可区分"；若完全一样（相同的球、相同的苹果），用"相同"。
2. \textbf{桶是否可区分？} 若桶有编号（不同的房间、不同的箱子），用"可区分"；若只是一堆无标号的盒子，用"相同"。
3. \textbf{是否允许空桶？} 允许空桶对应"非负整数解"，不允许空桶对应"正整数解"（相同球的情形）。

\textbf{例一。} 5 个不同的人住 3 个不同的房间，每间至少一人：这是模型 B，数目为

\[
3^5-3\cdot2^5+3=\sum_{i=0}^{3}(-1)^i\binom3i(3-i)^5=243-96+3=150.
\]

（也可以写成 $3!\,S(5,3)=6\cdot25=150$。）

\textbf{例二。} 5 个相同的苹果放进 3 个不同的篮子，允许空篮：模型 D，数目为 $\binom{5+3-1}{3-1}=\binom72=21$；若不允许空篮则为 $\binom{4}{2}=6$。

\textbf{例三。} 8 个不同的球放进 4 个不同的桶，允许空桶：模型 A，$4^8$；若每桶至少一个，则是 $\sum_{i=0}^{4}(-1)^i\binom4i(4-i)^8$。

\subsection{本节的勘误与提醒}


\begin{itemize}
\item \textbf{$m^n$ 还是 $n^m$？} $n$ 个球、$m$ 个桶时是 $m^n$：\textbf{每个球}有 $m$ 种选择。转写中"$nM$ 的 $n$ 次方"之类是口误，以 $m^n$ 为准。
\item \textbf{"补 $m$ 个球"的方向。} 手写稿写"若允许空桶，补 $m$ 个球"，意思是把"允许空"化成"每桶非空"：先给每桶补 1 个，再用非空公式；方向不要反。
\item \textbf{模型 B 的公式不能直接套模型 C 的公式。} 球可区分时不能用隔板法；满射必须用容斥（或第二类 Stirling 数乘 $m!$）。
\item \textbf{两种"相同"不要混。} "球相同、桶不同"是隔板法；"球相同、桶相同"是整数分拆，后者没有简单闭式。
\item 转写中"把 $n$ 个人分成 $k_1,\dots,k_t$ 组"与"球与桶"是两套说法：前者用多重组合数（有序分组），后者用 $m^n$ 或隔板法。
\end{itemize}

\section{容斥原理与抽屉原理}


这一节讲两个用途最广的"非构造性"计数/存在性工具：\textbf{容斥原理}（把"至少满足一个条件"的并集大小算清楚）与\textbf{抽屉原理}（保证某一类里至少有若干个对象）。容斥原理在第 4 节已经用来数过满射，这里补上一般形式与证明；抽屉原理则给出一个经典应用：Erdős–Szekeres 单调子列定理。

\textbf{本节用到的记号：} $[n]=\{1,\dots,n\}$；$A_i^c$ 为补集；$1_{P}$ 为谓词 $P$ 的指示函数（$P$ 真为 $1$，否则为 $0$）；$N$ 为全集大小；$D_n$ 为错排数；"子列"指保持原顺序选出的若干项。

\subsection{容斥原理}


\textbf{两个集合。} 对任意有限集 $A,B$，

\[
|A\cup B|=|A|+|B|-|A\cap B|.
\]

\textbf{三个集合。}

\[
|A\cup B\cup C|=|A|+|B|+|C|-|A\cap B|-|A\cap C|-|B\cap C|+|A\cap B\cap C|.
\]

\textbf{一般形式。} 设 $A_1,\dots,A_n$ 是有限集，则

\[
\left|\bigcup_{i=1}^{n}A_i\right|
=\sum_{\varnothing\neq S\subseteq[n]}(-1)^{|S|+1}\left|\bigcap_{i\in S}A_i\right|.
\]

即：单个集合大小之和，减去所有两两交集之和，加上所有三三交集之和，依次交替。

\textbf{证明一（按元素计数）。} 任取元素 $x$，设它恰好属于 $A_1,\dots,A_n$ 中的 $t$ 个集合。

\begin{itemize}
\item 若 $t=0$，$x$ 不在并集中，右边每一项都不含 $x$，贡献为 $0$；
\item 若 $t\ge1$，右边对 $x$ 的贡献是
\end{itemize}


\[
\sum_{s=1}^{t}(-1)^{s+1}\binom ts=1-\sum_{s=0}^{t}(-1)^{s}\binom ts=1-(1-1)^t=1,
\]

最后一步用二项式定理（$t\ge1$）。所以右边恰好把并集中的每个元素数了一次、并集外的元素数了零次。$\square$

\textbf{证明二（指示函数，课堂写法）。} 对每个 $x$ 定义 $b_i=1_{x\in A_i}\in\{0,1\}$。由二项式展开，

\[
\prod_{i=1}^{n}(1-b_i)=\sum_{S\subseteq[n]}(-1)^{|S|}\prod_{i\in S}b_i .
\]

若 $x$ 在并集中，则某个 $b_i=1$，左边为 $0$；于是

\[
1_{x\in\bigcup_i A_i}=1-\prod_{i=1}^{n}(1-b_i)=\sum_{\varnothing\neq S\subseteq[n]}(-1)^{|S|+1}\prod_{i\in S}b_i .
\]

而 $\prod_{i\in S}b_i=1_{x\in\bigcap_{i\in S}A_i}$。对 $x$ 求和即得一般形式。$\square$

\textbf{补集形式。} 设 $N$ 为全集大小，则

\[
|A_1^c\cap\cdots\cap A_n^c|
=N-\sum_i|A_i|+\sum_{i<j}|A_i\cap A_j|-\cdots+(-1)^n|A_1\cap\cdots\cap A_n|.
\]

它数的是"一个坏条件都不满足"的对象个数，是容斥最常见的用法。

\textbf{对称写法。} 若所有 $j$ 元交集的大小都相同（记为 $f(j)$，与具体指标无关），则并集公式简化为

\[
\left|\bigcup_{i=1}^{n}A_i\right|
=\sum_{j=1}^{n}(-1)^{j-1}\binom nj f(j).
\]

这在"各条件地位对称"的问题里很好用。例如数 $n$ 元排列中"至少有一个位置不动"的排列：取 $A_i$ 为"第 $i$ 个位置不动"的排列集合，则任意 $j$ 个 $A_i$ 的交集大小都是 $f(j)=(n-j)!$（与是哪 $j$ 个位置无关），于是

\[
\left|\bigcup_{i=1}^{n}A_i\right|=\sum_{j=1}^{n}(-1)^{j-1}\binom nj(n-j)!=n!-D_n,
\]

其中 $D_n$ 是错排数（与 §5.1.2 的公式一致）。

\textbf{例（$[100]$ 中与 $2,3,5$ 都互素的数）。} 记 $A_p$ 为 $p$ 的倍数集合，则 $|A_2|=50$、$|A_3|=33$、$|A_5|=20$；$|A_2\cap A_3|=16$（$6$ 的倍数）、$|A_2\cap A_5|=10$、$|A_3\cap A_5|=6$；$|A_2\cap A_3\cap A_5|=3$（$30$ 的倍数）。由补集形式，

\[
100-(50+33+20)+(16+10+6)-3=26.
\]

\subsubsection{应用一：满射计数}


设 $n,m\ge1$。令 $A_i$ 为"$[n]$ 到 $[m]$ 的函数中，$i$ 没有原像"的函数集合（即函数不取值为 $i$）。则

\[
|A_i|=(m-1)^n,\qquad
|A_{i_1}\cap\cdots\cap A_{i_k}|=(m-k)^n,
\]

因为共有 $k$ 个值被禁用，每个输入只剩 $m-k$ 个选择。满射是"一个值都不缺"的函数，所以

\[
\operatorname{onto}(n,m)=m^n-\left|\bigcup_{i=1}^{m}A_i\right|
=\sum_{k=0}^{m}(-1)^k\binom mk(m-k)^n .
\]

（$k=0$ 项为 $m^n$。）例：$n=5,m=3$ 时 $\operatorname{onto}(5,3)=3^5-3\cdot2^5+3\cdot1^5=243-96+3=150$，与第 4 节模型 B 一致。

\subsubsection{应用二：错排}


设 $A_i$ 为"把 $i$ 映到 $i$"的排列集合（$i$ 是"不动点"）。则

\[
|A_i|=(n-1)!,\qquad |A_{i_1}\cap\cdots\cap A_{i_k}|=(n-k)! .
\]

错排是"没有不动点"的排列，其个数为

\[
D_n=\sum_{k=0}^{n}(-1)^k\binom nk(n-k)!
=n!\sum_{k=0}^{n}\frac{(-1)^k}{k!}.
\]

前几项：$D_0=1$、$D_1=0$、$D_2=1$、$D_3=2$、$D_4=9$、$D_5=44$。

\subsubsection{应用三：Euler 函数（补充）}


设 $n=p_1^{a_1}\cdots p_r^{a_r}$，令 $A_i$ 为 $[n]$ 中被 $p_i$ 整除的数。$[n]$ 中与 $n$ 互素的数就是不在任何 $A_i$ 中的数。由容斥，单个交的大小为 $|A_i|=n/p_i$，$j$ 个交的大小为 $n/(p_{i_1}\cdots p_{i_j})$（与具体是哪 $j$ 个素数有关，故这里不用上面的对称写法），于是

\[
\begin{aligned}
\varphi(n)&=n-\sum_i\frac{n}{p_i}+\sum_{i<j}\frac{n}{p_ip_j}-\cdots+(-1)^r\frac{n}{p_1\cdots p_r}\\
&=n\prod_{i=1}^{r}\left(1-\frac1{p_i}\right).
\end{aligned}
\]

这是数论中 Euler 函数公式的容斥证明（教材 [DS] §11.3.5）。

\subsection{抽屉原理}


\textbf{基本形式。} 若把 $n+1$ 个对象放进 $n$ 个盒子，则至少有一个盒子里有 $\ge2$ 个对象。

\textbf{推广形式。} 若把 $kn+1$ 个对象放进 $n$ 个盒子（$k\ge1$），则至少有一个盒子里有 $\ge k+1$ 个对象。

\textbf{证明（反证）。} 若每个盒子至多 $k$ 个对象，则总数至多 $kn$，与 $kn+1$ 矛盾。$\square$

\textbf{例。} 任意 13 个人中必有两个人生日在同一个月（13 个对象、12 个盒子）；任意 $n+1$ 个整数中必有两个之差是 $n$ 的倍数（按模 $n$ 的余数分 $n$ 个盒子）。

\subsubsection{Erdős–Szekeres 定理}


\textbf{定理。} 设 $r,s\ge1$，$a_1,\dots,a_{rs+1}$ 是 $rs+1$ 个\textbf{两两不同}的实数。则其中一定存在长度至少 $r+1$ 的上升子列，或者长度至少 $s+1$ 的下降子列。

\textbf{证明。} 对每个位置 $i$，定义两个数：

\begin{itemize}
\item $u_i$：以 $a_i$ 结尾的最长上升子列的长度；
\item $v_i$：以 $a_i$ 结尾的最长下降子列的长度。
\end{itemize}


\textbf{关键观察：$(u_i,v_i)$ 两两不同。} 设 $i<j$。因为 $a_i\ne a_j$，有两种情况：

\begin{itemize}
\item 若 $a_i<a_j$：把 $a_j$ 接在"以 $a_i$ 结尾的最长上升子列"后面，得到长度 $u_i+1$ 的上升子列，所以 $u_j\ge u_i+1>u_i$；
\item 若 $a_i>a_j$：同理把 $a_j$ 接在最长下降子列后面，得到 $v_j\ge v_i+1>v_i$。
\end{itemize}


两种情况下 $(u_i,v_i)\ne(u_j,v_j)$。

\textbf{用抽屉原理推出矛盾。} 假设既没有长度 $r+1$ 的上升子列，也没有长度 $s+1$ 的下降子列，则所有 $u_i\le r$、$v_i\le s$，于是每个 $(u_i,v_i)$ 都落在大小为 $rs$ 的集合 $[r]\times[s]$ 中。但共有 $rs+1$ 个两两不同的数对，抽屉原理说明必有两个相同，矛盾。$\square$

\textbf{界是紧的。} 序列

\[
s,s-1,\dots,1,\quad 2s,2s-1,\dots,s+1,\quad\dots,\quad rs,rs-1,\dots,(r-1)s+1
\]

（共 $r$ 段、每段 $s$ 个，段内递减、段间递增）长度为 $rs$，上升子列最长恰为 $r$（每段至多取一个），下降子列最长恰为 $s$（必须在一段内），所以 $rs$ 个数确实可以既没有长 $r+1$ 上升、也没有长 $s+1$ 下降。这说明定理中的 $rs+1$ 不能改成 $rs$。

\textbf{关于"不同"的条件。} 定理要求元素两两不同，否则全是相等元素时没有长度 $2$ 的严格上升或下降子列。若元素可以重复，应把结论改成"非降子列或非升子列"（教材版本），证明中把"$a_i<a_j$"换成"$a_i\le a_j$"并注意相等时的处理即可。

\subsubsection{两个补充例子（教材 §11.4）}


\begin{itemize}
\item \textbf{互素对。} 若 $S\subseteq[2n]$ 且 $|S|\ge n+1$，则 $S$ 中存在两个互素的数。证明：按 $\{1,2\},\{3,4\},\dots,\{2n-1,2n\}$ 分成 $n$ 个"抽屉"，每个抽屉里的两个数互素；由抽屉原理，某个抽屉的两个数都在 $S$ 中。
\item \textbf{整除对。} 同样的 $S$ 中还存在 $i\ne j$ 使 $i\mid j$。证明：把每个数写成"奇数部分 $\times\,2^k$"，按奇数部分（共 $n$ 个可能的奇数）分抽屉；同抽屉的两个数一个是另一个的 $2$ 的幂倍，故有整除关系。
\end{itemize}


\subsection{本节的勘误与提醒}


\begin{itemize}
\item \textbf{容斥的符号。} 一般形式是 $(-1)^{|S|+1}$：一个集合取 $+$，两个集合取 $-$，三个集合取 $+$，依次交替。补集形式则从 $N$ 开始、第一项取 $-$，两种写法不要混用。
\item \textbf{满射公式的 $i=0$ 项。} 是 $m^n$（不是 $0$）；$n<m$ 时公式自动给出 $0$。
\item \textbf{错排的约定。} $D_0=1$（空排列没有不动点，算一种）；$D_1=0$。
\item \textbf{抽屉原理的"$kn+1$"。} 是 $kn+1$ 个对象放进 $n$ 个盒子才能保证某盒有 $k+1$ 个；$kn$ 个只能保证某盒有 $k$ 个。
\item \textbf{Erdős–Szekeres 的长度是"至少 $r+1$ / $s+1$"。} 转写里说"长度超过 $r$ 的上升子列或长度超过 $s$ 的下降子列"，与"至少 $r+1$ / $s+1$"等价。
\item \textbf{严格与非严格。} 有重复元素时不能直接套"严格上升/下降"的版本；要么先假设元素互异，要么改用"非降/非升"的教材版本。
\item 课堂转写里"抽屉原理的第几首"是"第几个例子"的口误；以本节内容为准。
\end{itemize}


\section*{一页式总览}
\begin{itemize}
\item 加法原理：不交并的大小相加；乘法原理：$|A\times B|=|A||B|$，分步独立相乘。
\item $P(n,k)=n^{\underline{k}}=\frac{n!}{(n-k)!}$（单射计数）；$\binom nk=\frac{n^{\underline{k}}}{k!}$（子集计数）；$\binom{n}{k_1,\dots,k_t}=\frac{n!}{k_1!\cdots k_t!}$（有序分组）。
\item 基本恒等式：$\binom nk=\binom n{n-k}$；$\binom nk=\binom{n-1}{k-1}+\binom{n-1}k$；$k\binom nk=n\binom{n-1}{k-1}$；$\binom nm\binom mk=\binom nk\binom{n-k}{m-k}$。
\item 二项式定理：$(x+y)^n=\sum_k\binom nk x^ky^{n-k}$；$\sum_k\binom nk=2^n$；$\sum_k(-1)^k\binom nk=0$；多项式定理：$(x_1+\cdots+x_t)^n=\sum_{k_1+\cdots+k_t=n}\binom{n}{k_1,\dots,k_t}x_1^{k_1}\cdots x_t^{k_t}$。
\item 曲棍球：$\sum_{k=m}^n\binom km=\binom{n+1}{m+1}$；一般斜向求和 $\sum_{k=0}^n\binom{r+k}{k}=\binom{r+n+1}{n}$（$r=n$ 时为 $\binom{2n+1}{n}$）；Vandermonde：$\sum_k\binom rk\binom m{n-k}=\binom{r+m}{n}$，特例 $\sum_k\binom nk^2=\binom{2n}{n}$。
\item 加权和：$\sum_k\binom kj\binom nk=\binom nj2^{n-j}$，于是 $\sum_k k\binom nk=n2^{n-1}$、$\sum_k k^2\binom nk=n(n+1)2^{n-2}$；任何 $k$ 的多项式都可化成 $\binom kj$ 的线性组合再求和。
\item 广义二项式系数：$\binom nk=\frac{n^{\underline k}}{k!}$（$n$ 可为负整数），$\binom{-1}{k}=(-1)^k$，$\binom{-m}{k}=(-1)^k\binom{m+k-1}{k}$；帕斯卡与吸收恒等式仍成立。\item 注意两种容易混淆的求和：$A_m=\sum_{k\ge0}\binom{k}{m}2^{-k}\equiv2$（手写稿底部），而 $A_n=\sum_{k=0}^n\binom{n}{k}2^{-k}=(3/2)^n$；两者递推不同。
\item $A_m=\sum_k(-1)^k\binom{m-k}{k}$ 满足 $A_m=A_{m-1}-A_{m-2}$，$A_0=A_1=1$，周期为 $6$（$1,1,0,-1,-1,0,\dots$）；二阶线性递推用特征方程 $\lambda^2=\alpha\lambda+\beta$ 求通解（相异根 $\lambda_1^n,\lambda_2^n$；重根 $(C_1+C_2m)\lambda^m$）。
\item 卡特兰数：$C_n=\frac1{n+1}\binom{2n}{n}$，递推 $C_n=\sum_{k=0}^{n-1}C_kC_{n-1-k}$；计数 $n$ 对括号、Dyck 格路、$(n+2)$ 边形三角剖分、$n+1$ 个符号的加括号方式。
\item 球与桶：$m^n$（可区分球/桶、允许空）；$\sum_i(-1)^i\binom mi(m-i)^n$（可区分球/桶、非空=满射）；$\binom{n-1}{m-1}$（相同球、非空）；$\binom{n+m-1}{m-1}$（相同球、允许空）。
\item 容斥：$|\bigcup_iA_i|=\sum_{\varnothing\ne S\subseteq[n]}(-1)^{|S|+1}|\bigcap_{i\in S}A_i|$；满射数 $=\sum_i(-1)^i\binom mi(m-i)^n$；错排 $D_n=n!\sum_{k=0}^n(-1)^k/k!$。
\item 抽屉原理：$kn+1$ 个对象放入 $n$ 个盒子 $\Rightarrow$ 某盒 $\ge k+1$ 个；Erdős--Szekeres：$rs+1$ 个两两不同的数必有长 $r+1$ 上升子列或长 $s+1$ 下降子列。
\end{itemize}

\section*{参考资料}
\begin{itemize}
\item \textbf{[DS]} 邓小铁、王畅。《离散数学与结构》§11.1（基本计数模型）、§11.2（组合恒等式）、§11.3（容斥原理）、§11.4（抽屉原理与 Erdős--Szekeres）。§13（生成函数）属于后续内容。
\item \textbf{[讲义]} 课程讲义 \texttt{25-26-lecture\_notes.pdf} 第九章：§9.1 基本计数原理、§9.2 组合数性质、§9.3 二项式定理的例子、§9.4 卡特兰数、§9.5 Balls \& Bins、§9.6 多项式推广、§9.7 容斥原理、§9.8 鸽笼原理；§9.9 生成函数为下一讲内容。
\item 原始材料：\texttt{DiscreteCourse2026/课堂手记和录音转写/Lec12/}（两段转写与三页手写原稿）。
\end{itemize}
\end{document}
